
\section{Additional Rules}

Last week I mentioned that we will add some of the valid arguments we prove
to our rule system. We start by defining the ones we proved last week.

\begin{fitchrule}[Disjunctive Syllogism]
  \label{def:disjunctive-syllogism-rule}
  \leavevmode
  % \vspace{-0.9\baselineskip}
  There are two rules for disjunctive syllogism, one for each disjunct:\\
  \twoitems{  
    \begin{fitchproof}
      \have[m] {1} {\ephi\eor\epsi}
      \have[n] {2} {\enot\ephi}
      \have[ ] {3} {\epsi}          \by{DS}{1,2}
    \end{fitchproof}
  }{
    \begin{fitchproof}
      \have[m] {1} {\ephi\eor\epsi}
      \have[n] {2} {\enot\epsi}
      \have[ ] {3} {\ephi}          \by{DS}{1,2}
    \end{fitchproof}
  }
\end{fitchrule}

\begin{fitchrule}[Excluded Middle]
  \label{def:excluded-middle-rule}
  \leavevmode
  \vspace{-0.9\baselineskip}
  \begin{fitchproof}
    \open
    \hypo[i] {1} {\ephi}      \by{AS}{}
    \have[j] {2} {\epsi}
    \close
    \open
    \hypo[k] {3} {\enot\ephi} \by{AS}{}
    \have[l] {4} {\epsi}
    \close
    \have[ ] {5} {\epsi}      \by{LEM}{1-2,3-4}
  \end{fitchproof}    
\end{fitchrule}

A few more rules will be extremely helpful. The first is that we can treat the
negation of a negation as no negation at all.

\begin{fitchrule}[Double Negative Elimination]
  \label{def:double-negative-rule}
  \leavevmode
  \vspace{-0.9\baselineskip}
  \begin{fitchproof}
    \have[m] {1} {\enot\enot\ephi}
    \have[ ] {2} {\ephi} \by{DNE}{1}
  \end{fitchproof}
\end{fitchrule}

The proof for this is simple.

\begin{fitchproof}
  \have {1} {\enot\enot\ephi} \by{PR}{}
  \open
  \hypo {2} {\enot\ephi}      \by{AS}{}
  \have {3} {\ebot}           \ne{1,2}
  \close
  \have {4} {\ephi}           \by{IP}{2-3}
\end{fitchproof}

Another useful technique, called \emph{modus tollens} lets us conclude that
the antecedant of a true conditional is false when its consequent is. The
rule works much like disjunctive syllogism.

\begin{fitchrule}[Modus Tollens]
  \label{def:modus-tollens-rule}
  \leavevmode
  \vspace{-0.9\baselineskip}
  \begin{fitchproof}
    \have[m] {1} {\ephi\eif\epsi}
    \have[n] {2} {\enot\epsi}
    \have[ ] {3} {\enot\ephi}     \by{MT}{1}
  \end{fitchproof}
\end{fitchrule}

This is another proof that leverages \ebot.

\begin{fitchproof}
  \have {1} {\ephi\eif\epsi}  \by{PR}{}
  \have {2} {\enot\epsi}      \by{PR}{}
  \open
  \hypo {3} {\ephi}           \by{AS}{}
  \have {4} {\epsi}           \ie{1,3}
  \have {5} {\ebot}           \ne{2,4}
  \close
  \have {6} {\enot\ephi}      \ni{3-5}
\end{fitchproof}

Finally, the \emph{De Morgan rules} tell us how to properly ditribute negation
into conjunctions and disjunctions.

\begin{fitchrule}[De Morgan Rules]
  \label{def:de-morgan-rules}
  \leavevmode
  % \vspace{-0.9\baselineskip}
  There are four rules for disjunctive syllogism forming a square of duals
  and rule reversals.

  We have a pair of rules for negated conjunctions: \\
  \twoitems{  
    \begin{fitchproof}
      \have[m] {1} {\enot(\ephi\eand\epsi)}
      \have[ ] {2} {\enot\ephi\eor\enot\epsi}   \by{DS}{1}
    \end{fitchproof}
  }{
    \begin{fitchproof}
      \have[m] {1} {\enot\ephi\eor\enot\epsi}
      \have[ ] {2} {\enot(\ephi\eand\epsi)}     \by{DS}{1}
    \end{fitchproof}
  }

  And a dual pair of rules for negated disjunctions: \\
  \twoitems{  
    \begin{fitchproof}
      \have[m] {1} {\enot(\ephi\eor\epsi)}
      \have[ ] {2} {\enot\ephi\eand\enot\epsi}  \by{DS}{1}
    \end{fitchproof}
  }{
    \begin{fitchproof}
      \have[m] {1} {\enot\ephi\eand\enot\epsi}
      \have[ ] {2} {\enot(\ephi\eor\epsi)}      \by{DS}{1}
    \end{fitchproof}
  }
\end{fitchrule}

We will look at these proofs next time, when we work through questions about
the exercises.

% TODO: fold the De Morgan proofs into an exercises document

